
\subsubsection{6.1 Interpretation}

The results indicate that hidden states at layer 25 of Llama-3-8B-Instruct, extracted at the TBG position, do not encode semantic entropy information accessible to logistic regression probes under the tested configuration.

Three competing explanations for this failure are:

\textbf{Layer localization:} SE-relevant information may be concentrated in layers other than layer 25. Published work reports optimal configurations without characterizing failure modes of suboptimal choices.

\textbf{Token position sensitivity:} The TBG position captures the model state before generation. SE information may emerge during generation and be better captured at alternative positions.

\textbf{Implementation factors:} Accumulated differences in response generation, SE computation, or preprocessing relative to published implementations may contribute to the performance gap.

The experiment cannot distinguish among these explanations.

\subsubsection{6.2 Similarity Augmentation}

The hypothesis that similarity features would improve SE prediction was directionally supported (positive delta) but the effect was negligible. Contributing factors may include:

\begin{enumerate}
\item \textbf{Feature dimensionality imbalance:} 4 similarity features relative to 4096 hidden state dimensions limits the influence of similarity information.
\end{enumerate}

\begin{enumerate}
\item \textbf{Absence of base signal:} When hidden states contain no extractable SE information, auxiliary features cannot compensate.
\end{enumerate}

\begin{enumerate}
\item \textbf{Linear probe limitations:} Logistic regression may not capture complex interactions between hidden states and similarity features.
\end{enumerate}

\subsubsection{6.3 Limitations}

This work has several limitations that constrain the conclusions:

\textbf{Single layer tested:} Only layer 25 of 32 was evaluated. SE information may exist at other layers.

\textbf{Single token position:} Only TBG was tested. Alternative positions (SLT, pooled representations) were not explored.

\textbf{Single dataset:} TruthfulQA is adversarial by design; results may not generalize to other benchmarks.

\textbf{Linear probe architecture:} Logistic regression may be insufficient; MLP probes may capture patterns missed by linear models.

\textbf{SE label quality not validated:} DeBERTa NLI clustering was assumed to produce valid SE labels without explicit quality verification.

\textbf{Single random seed:} No variance estimation across multiple seeds.

These limitations are acknowledged as acceptable for an existence proof: the experiment tested whether a reasonable configuration works, and it did not. Systematic exploration remains for future work.

\subsubsection{6.4 Implications}

For practitioners considering SE probe deployment:

\begin{enumerate}
\item Published configurations may not transfer across implementations or model versions.
\item Systematic validation across layers, token positions, and probe architectures may be necessary.
\item AUROC close to 0.50 on held-out data should trigger configuration review before deployment.
\end{enumerate}

\subsubsection{6.5 Future Directions}

The negative result suggests several directions for investigation:

\begin{enumerate}
\item \textbf{Layer ablation:} Test layers 20-31 to locate where, if anywhere, SE information resides in Llama-3-8B-Instruct.
\end{enumerate}

\begin{enumerate}
\item \textbf{Token position comparison:} Compare TBG, SLT, and pooled representations.
\end{enumerate}

\begin{enumerate}
\item \textbf{Nonlinear probes:} Test MLP architectures with hidden layers.
\end{enumerate}

\begin{enumerate}
\item \textbf{SE label validation:} Inspect semantic clustering quality before probe training.
\end{enumerate}

\begin{enumerate}
\item \textbf{Multi-dataset evaluation:} Validate on TriviaQA and SQuAD to assess generalization.
\end{enumerate}
